% Research Talk: a 16:9 Beamer deck with the hand-built LabShelf theme.
% The theme (colours, fonts, footline, divider and statement slides) lives in
% beamerthemelabshelftalk.sty; this file holds only the talk itself.
\documentclass[aspectratio=169,10pt]{beamer}

\usetheme{labshelftalk}
\usepackage{booktabs}
\usepackage[style=authoryear,maxcitenames=2,maxbibnames=3,uniquelist=false,backend=biber]{biblatex}
\addbibresource{refs.bib}
\renewcommand*{\bibfont}{\scriptsize}

% Title metadata: the short forms appear in the footline of every slide.
\title{Cool streets}
\subtitle{Three decades of urban heat islands in twelve European cities, seen from orbit}
\author[M. Okafor]{Maya Okafor\texorpdfstring{\quad}{, }Tom\'as Rivera\texorpdfstring{\quad}{, }Lena Hartmann}
\institute{Institute for Urban Climate, Northfield University}
\date{Urban Climate Seminar \enspace\textperiodcentered\enspace 14 May 2026}

\begin{document}

\begin{frame}[plain]
  \titlepage
\end{frame}

\section{Motivation}

\begin{frame}{Heat is Europe's deadliest weather hazard}
  \vspace{5mm}
  \begin{columns}[T,onlytextwidth]
    \column{0.3\textwidth}
      {\lsdisplay{28}\color{lsAccent}61\,672\par}
      \vspace{3mm}\rule{\linewidth}{0.4pt}\par\vspace{2mm}
      {\small heat-related deaths in Europe in the summer of 2022 \parencite{ballester2023}\par}
    \column{0.3\textwidth}
      {\lsdisplay{28}\color{lsAccent}+4.6\,K\par}
      \vspace{3mm}\rule{\linewidth}{0.4pt}\par\vspace{2mm}
      {\small mean summer surface heat island of the twelve cities we study, 2024\par}
    \column{0.3\textwidth}
      {\lsdisplay{28}\color{lsAccent}1 in 3\par}
      \vspace{3mm}\rule{\linewidth}{0.4pt}\par\vspace{2mm}
      {\small Europeans live in a city of more than 100\,000 people\par}
  \end{columns}
  \vspace{0pt plus 1filll}
  {\footnotesize\color{lsMuted}Built surfaces store heat by day and release it slowly
  \parencite{oke1982}; the street, not the weather station, is where people feel it.\par}
  \vspace{4mm}
\end{frame}

\section{Data and method}

\begin{frame}{Thirty-four summers of thermal imagery}
  \begin{columns}[T,onlytextwidth]
    \column{0.54\textwidth}
      \begin{itemize}
        \item 1\,264 cloud-free summer scenes from Landsat 5, 7, 8 and 9
        \item Twelve cities, from Lisbon to Warsaw, each with a rural reference ring
        \item Thermal bands aggregated from 30\,m to a 100\,m street grid
        \item 4\,812 neighbourhoods linked to canopy, imperviousness and building height
      \end{itemize}
      \vspace{3mm}
      {\footnotesize\color{lsMuted}Scenes acquired between June and August, about 10:30 local time.\par}
    \column{0.38\textwidth}
      \centering
      \begin{tikzpicture}
        \lsheatraster{12}{12}{2.8mm}{lsPaper}
        \foreach \c [count=\k from 0] in {0,20,40,60,80,100}{
          \fill[lsAccent!\c!lsPaper] (\k*5.6mm,-3.6mm) rectangle ++(4.8mm,1.4mm);}
        \node[anchor=north west,inner sep=0,font=\tiny\ttfamily,text=lsMuted] at (0,-5mm) {COOLER};
        \node[anchor=north east,inner sep=0,font=\tiny\ttfamily,text=lsMuted] at (32.8mm,-5mm) {HOTTER};
      \end{tikzpicture}\par
      \vspace{2mm}
      {\scriptsize\color{lsMuted}Land surface temperature, Madrid centre, 3 August 2024\par}
  \end{columns}
\end{frame}

\begin{frame}{From radiance to surface temperature}
  \vspace{9mm}
  {\Large\begin{equation*}
    T_s \;=\; \frac{\tikzmarknode{tb}{T_B}}
      {1 + \bigl(\tikzmarknode{lam}{\lambda}\, T_B / \tikzmarknode{rho}{\rho}\bigr)\,
       \ln \tikzmarknode{eps}{\varepsilon}}
  \end{equation*}}
  \begin{tikzpicture}[remember picture,overlay,
      note/.style={font=\scriptsize,text=lsInk,align=left,inner sep=0},
      lead/.style={lsAccent,line width=0.5pt}]
    \draw[lead] (tb.north) -- ++(0,4mm) -- ++(-20mm,0) coordinate (a);
    \node[note,anchor=east,align=right] at ($(a)+(-1.5mm,0)$) {{\lssemibold at-sensor brightness}\\ temperature, Band 10};
    \draw[lead] (lam.south) -- ++(0,-7mm) -- ++(-16mm,0) coordinate (b);
    \node[note,anchor=east,align=right] at ($(b)+(-1.5mm,0)$) {{\lssemibold wavelength}\\ 10.9\,\textmu m};
    \draw[lead] (rho.south) -- ++(0,-7mm) -- ++(16mm,0) coordinate (c);
    \node[note,anchor=west] at ($(c)+(1.5mm,0)$) {{\lssemibold radiation constant}\\ $hc/k_B = 1.438$\,cm\,K};
    \draw[lead] (eps.north) -- ++(0,11mm) -- ++(14mm,0) coordinate (d);
    \node[note,anchor=west] at ($(d)+(1.5mm,0)$) {{\lssemibold surface emissivity}\\ from NDVI};
  \end{tikzpicture}\par
  \vspace{0pt plus 1filll}
  {\footnotesize\color{lsMuted}Single-channel retrieval after \textcite{jimenez2003}. Heat-island
  intensity is the urban mean minus the rural-ring mean, scene by scene.\par}
  \vspace{4mm}
\end{frame}

\section{Results}

\begin{frame}{Heat islands have grown in every city we studied}
  \begin{columns}[T,onlytextwidth]
    \column{0.68\textwidth}
      \begin{tikzpicture}
        \begin{axis}[lsaxis, width=\linewidth-24mm, height=36mm,
            xmin=1990, xmax=2024, ymin=2, ymax=6,
            xtick={1990,2000,2010,2020}, ytick={2,3,4,5,6},
            x tick label style={/pgf/number format/1000 sep={}},
            ylabel={Surface heat island (K)}]
          \addplot[lsRule!40!lsMuted] coordinates {(1990,2.4)(1992,2.5)(1994,2.4)(1996,2.6)(1998,2.6)(2000,2.7)(2002,2.8)(2004,2.8)(2006,2.9)(2008,3.0)(2010,3.0)(2012,3.1)(2014,3.2)(2016,3.2)(2018,3.3)(2020,3.4)(2022,3.5)(2024,3.6)}
            node[right,font=\scriptsize,text=lsMuted] {Lisbon};
          \addplot[lsRule!40!lsMuted] coordinates {(1990,2.8)(1992,2.9)(1994,3.0)(1996,2.9)(1998,3.2)(2000,3.3)(2002,3.3)(2004,3.5)(2006,3.6)(2008,3.7)(2010,3.8)(2012,3.8)(2014,4.0)(2016,4.1)(2018,4.3)(2020,4.3)(2022,4.5)(2024,4.6)}
            node[right,font=\scriptsize,text=lsMuted] {Athens};
          \addplot[lsAccent,line width=1.6pt] coordinates {(1990,3.1)(1992,3.3)(1994,3.2)(1996,3.5)(1998,3.6)(2000,3.6)(2002,3.9)(2004,4.0)(2006,4.1)(2008,4.3)(2010,4.2)(2012,4.5)(2014,4.7)(2016,4.6)(2018,4.9)(2020,5.1)(2022,5.2)(2024,5.4)}
            node[right,font=\scriptsize\lssemibold,text=lsAccent] {Madrid};
        \end{axis}
      \end{tikzpicture}
    \column{0.26\textwidth}
      \vspace{2mm}
      {\lsdisplay{22}\color{lsAccent}+0.7\,K\par}
      \vspace{1.5mm}
      {\footnotesize per decade in Madrid, the fastest of the twelve.\par}
      \vspace{5mm}
      {\footnotesize\color{lsMuted}Mean of all summer scenes per two-year window. All twelve
      trends are positive at $p<0.01$, in line with the global survey of \textcite{peng2012}.\par}
  \end{columns}
\end{frame}

\begin{frame}{Tree canopy is the strongest cooling lever}
  \begin{columns}[T,onlytextwidth]
    \column{0.68\textwidth}
      \begin{tikzpicture}
        \begin{axis}[lsaxis, width=\linewidth-24mm, height=36mm,
            xmin=0, xmax=48, ymin=-5, ymax=4,
            xtick={0,10,20,30,40}, ytick={-4,-2,0,2,4},
            xlabel={Tree canopy cover (\%)},
            ylabel={Temperature anomaly (K)}]
          \addplot[only marks,mark=*,mark size=1.3pt,lsAccent,fill opacity=0.55,draw opacity=0]
            coordinates {(16.2,0.78)(5.2,3.06)(25.6,-1.59)(24.3,-1.08)(3.6,2.52)(6.0,2.35)
            (20.7,-0.40)(11.8,0.90)(29.6,-1.54)(19.5,-0.53)(45.0,-4.10)(14.7,0.91)
            (8.3,2.16)(37.9,-3.48)(10.0,0.83)(18.4,-0.48)(26.1,-1.38)(11.1,1.37)
            (31.9,-3.10)(27.8,-1.67)(21.9,-1.07)(32.8,-1.89)(12.7,0.38)(40.5,-4.60)
            (34.1,-3.44)(7.2,3.53)(20.4,-0.46)(23.5,-1.38)(3.7,2.23)(27.2,-2.59)
            (40.5,-4.63)(28.2,-1.17)(27.5,-2.84)(43.6,-4.59)(22.9,-1.04)(32.9,-3.01)
            (30.5,-1.37)(14.5,0.60)(19.0,-0.26)(22.3,-0.94)(9.4,1.76)(35.8,-3.27)
            (7.7,1.95)(40.3,-3.72)(5.5,1.69)(40.9,-4.15)};
          \addplot[lsInk,line width=0.8pt,domain=2:46,samples=2] {3.4-0.19*x};
        \end{axis}
      \end{tikzpicture}
    \column{0.26\textwidth}
      \vspace{2mm}
      {\lsdisplay{22}\color{lsAccent}$-$1.9\,K\par}
      \vspace{1.5mm}
      {\footnotesize for every ten points of canopy cover.\par}
      \vspace{5mm}
      {\footnotesize\color{lsMuted}Each dot is a 100-neighbourhood bin from all twelve cities;
      the line is the pooled fit. Consistent with \textcite{schwaab2021}.\par}
  \end{columns}
\end{frame}

\begin{frame}{What drives the street-level anomaly}
  \vspace{3mm}
  \centering\small\renewcommand{\arraystretch}{1.3}
  \begin{tabular}{@{}l r r r@{}}
    \toprule
    {\lssemibold Driver} & {\lssemibold Effect (K)} & {\lssemibold 95\% interval} & {\lssemibold Variance} \\
    \midrule
    Impervious surface, +10 pp   & $+1.42$ & $[\,1.31,\ 1.53\,]$ & 38\% \\
    Tree canopy, +10 pp          & \textcolor{lsAccent}{$-1.87$} & $[-2.01,\ -1.73\,]$ & 29\% \\
    Roof albedo, +0.1            & $-0.94$ & $[-1.10,\ -0.78\,]$ & 12\% \\
    Mean building height, +5\,m  & $+0.36$ & $[\,0.22,\ 0.50\,]$ & 9\% \\
    Distance to water, +1\,km    & $+0.21$ & $[\,0.12,\ 0.30\,]$ & 5\% \\
    \bottomrule
  \end{tabular}\par
  \vspace{4mm}
  {\footnotesize\color{lsMuted}Mixed-effects model with a random intercept per city;
  $n = 4\,812$ neighbourhoods, $R^2 = 0.71$.\par}
\end{frame}

\lsstatement{Key takeaway}
  {Ten more points of tree canopy cool a street by \textcolor{lsAccent}{almost two degrees}, more than any other lever a city controls.}
  {Pooled across twelve cities and thirty-four summers of Landsat thermal imagery.}

\section{Implications}

\begin{frame}{Three levers for cooler streets}
  \vspace{4mm}
  \begin{columns}[T,onlytextwidth]
    \column{0.3\textwidth}
      {\ttfamily\small\color{lsAccent}01\par}\vspace{1mm}\rule{\linewidth}{0.4pt}\par\vspace{2mm}
      {\lssemibold Plant where it is hottest\par}\vspace{1mm}
      {\footnotesize Target canopy at the 10\% of blocks with the highest anomaly; they
      hold 31\% of residents over 65.\par}
    \column{0.3\textwidth}
      {\ttfamily\small\color{lsAccent}02\par}\vspace{1mm}\rule{\linewidth}{0.4pt}\par\vspace{2mm}
      {\lssemibold Lighten the roofs\par}\vspace{1mm}
      {\footnotesize Raising roof albedo by 0.1 recovers about half the effect of new
      canopy, at a fraction of the cost.\par}
    \column{0.3\textwidth}
      {\ttfamily\small\color{lsAccent}03\par}\vspace{1mm}\rule{\linewidth}{0.4pt}\par\vspace{2mm}
      {\lssemibold Measure from orbit\par}\vspace{1mm}
      {\footnotesize The pipeline runs on public Landsat data, so any city can track its
      heat island every summer.\par}
  \end{columns}
\end{frame}

\section*{References}

\begin{frame}{References}
  \printbibliography[heading=none]
\end{frame}

\begin{frame}[plain]
  \begin{tikzpicture}[remember picture,overlay]
    \begin{scope}[shift={($(current page.north east)+(-72mm,-14mm-69mm)$)}]
      \lsheatraster{15}{15}{4mm}{lsPaper}
    \end{scope}
    \node[anchor=north west,inner sep=0] at ($(current page.north west)+(14mm,-14mm)$)
      {\lskicker{Questions}};
    \node[anchor=west,inner sep=0] at ($(current page.west)+(14mm,4mm)$)
      {\fontsize{40}{40}\lslight\selectfont\color{lsInk}Thank you};
    \fill[lsAccent] ($(current page.south west)+(14mm,23mm)$) rectangle ++(16mm,0.8pt);
    \node[anchor=north west,inner sep=0,text width=80mm]
      at ($(current page.south west)+(14mm,19mm)$) {%
      {\small\lssemibold maya.okafor@northfield.edu\par}
      \vspace{1mm}{\scriptsize\ttfamily\color{lsMuted}data and code \enspace github.com/northfield-climate/cool-streets\par}};
  \end{tikzpicture}
\end{frame}

\appendix

\begin{frame}{Backup: robust to the emissivity scheme}
  \begin{tikzpicture}
    \begin{axis}[lsaxis, xbar, width=0.7\textwidth, height=30mm, scale only axis,
        bar width=3.2mm, bar shift=0pt, xmin=0, xmax=2.4, enlarge y limits=0.18,
        grid=none, xmajorgrids,
        symbolic y coords={Fixed 0.97,Land-cover lookup,NDVI threshold,NDVI mixture},
        ytick={Fixed 0.97,Land-cover lookup,NDVI threshold,NDVI mixture}, y tick label style={font=\scriptsize,text=lsInk},
        xlabel={Cooling per +10 pp canopy (K)},
        nodes near coords, nodes near coords align={horizontal},
        every node near coord/.append style={font=\scriptsize\ttfamily,text=lsMuted}]
      \addplot[fill=lsRule!70!lsMuted,draw=none] coordinates
        {(1.62,Fixed 0.97) (1.79,Land-cover lookup) (1.84,NDVI threshold)};
      \addplot[fill=lsAccent,draw=none] coordinates {(1.87,NDVI mixture)};
    \end{axis}
  \end{tikzpicture}\par
  \vspace{2mm}
  {\footnotesize\color{lsMuted}The main analysis uses the NDVI mixture scheme (highlighted).
  Every alternative keeps canopy as the strongest cooling lever.\par}
\end{frame}

\end{document}
